\begin{pastexam}[label={px:df3-106-ds-04}]{地方三等資訊處理 106 年 第四題}
（申論；20 分）

給定下列數列，若以快速排序法（Quick Sort）、選擇排序法（Selection Sort）、堆積排序法（Heap Sort）、
泡沫排序法（Bubble Sort）進行排序。請問下列數列是那一個排序法排序過程的暫時結果，並說明之。
（每小題 5 分，共 20 分）
\begin{center}
\small
\begin{tabular}{r|*{16}{c|}}
\cline{2-17}
 & 75 & 93 & 32 & 81 & 75 & 89 & 89 & 99 & 25 & 78 & 54 & 75 & 87 & 12 & 75 & 28 \\
\cline{2-17}
\multicolumn{17}{c}{} \\[-0.6em]
\cline{2-17}
㈠ & 99 & 93 & 89 & 81 & 78 & 87 & 89 & 75 & 25 & 75 & 54 & 75 & 32 & 12 & 75 & 28 \\
\cline{2-17}
\multicolumn{17}{c}{} \\[-0.6em]
\cline{2-17}
㈡ & 25 & 28 & 32 & 75 & 12 & 75 & 54 & 75 & 99 & 78 & 89 & 89 & 87 & 75 & 81 & 93 \\
\cline{2-17}
\multicolumn{17}{c}{} \\[-0.6em]
\cline{2-17}
㈢ & 12 & 25 & 28 & 32 & 54 & 75 & 75 & 75 & 75 & 78 & 81 & 87 & 89 & 99 & 93 & 89 \\
\cline{2-17}
\multicolumn{17}{c}{} \\[-0.6em]
\cline{2-17}
㈣ & 32 & 75 & 75 & 81 & 89 & 25 & 78 & 54 & 75 & 87 & 12 & 75 & 28 & 89 & 93 & 99 \\
\cline{2-17}
\end{tabular}
\end{center}
\end{pastexam}

\begin{pastsoln}
四個數列各對應一種方法，看的是各方法排到一半時「留下的痕跡」（資結書第~14~章）。 原數列共 16 個數，
排好是 12, 25, 28, 32, 54, 75, 75, 75, 75, 78, 81, 87, 89, 89, 93, 99。 位置編號從 1 開始。

\textbf{㈠ 堆積排序：建好最大堆積之後的樣子。} 檢查它是不是最大堆積（$A[i] \ge A[2i]$、$A[2i+1]$）：
99 $\ge$ 93, 89；93 $\ge$ 81, 78；89 $\ge$ 87, 89；81 $\ge$ 75, 25；78 $\ge$ 75, 54；87 $\ge$ 75, 32；
89 $\ge$ 12, 75；75 $\ge$ 28。 全部成立。 實際對原數列做 Build-Max-Heap（從 $i = 8$ 往回到 $i = 1$
逐一往下調整，資結書第~8.2~節），結果正好就是這個數列。 堆積排序的第一階段就是建堆積，接下來才
反覆把根（最大值）換到最後面，所以這是堆積排序建堆完成、還沒開始取出時的暫時結果。 它不是另外三種：
最大的 99 在最前面、整體大致由大到小，由小到大排序的其他方法不會出現這種樣子。

\textbf{㈡ 快速排序：第一次分割（partition）之後。} 以第一個數 75 為 pivot，比 75 小或等於的放左邊、
大或等於的放右邊，最後把 pivot 放到中間。 數列裡第 8 個位置是 75，它左邊 25, 28, 32, 75, 12, 75, 54
都 $\le 75$，右邊 99, 78, 89, 89, 87, 75, 81, 93 都 $\ge 75$，左右兩邊各自還沒排序：這正是分割完的樣子。
用課本常見的分割法（Horowitz 版：$i$ 從左往右找 $\ge$ pivot 的，$j$ 從右往左找 $\le$ pivot 的，交換，
直到 $i$、$j$ 交錯，再把 pivot 和 $A[j]$ 交換）實際做一次：
\begin{center}
\begin{tabular}{ccl}
\toprule
$i$ 停在 & $j$ 停在 & 動作 \\
\midrule
2（93） & 16（28） & 交換 \\
4（81） & 15（75） & 交換 \\
5（75） & 14（12） & 交換 \\
6（89） & 12（75） & 交換 \\
7（89） & 11（54） & 交換 \\
8（99） & 9（25） & 交換 \\
9 & 8 & 交錯，pivot 和 $A[8]$ 交換 \\
\bottomrule
\end{tabular}
\end{center}
得到 25, 28, 32, 75, 12, 75, 54, \textbf{75}, 99, 78, 89, 89, 87, 75, 81, 93，和數列完全一樣。

\textbf{㈢ 選擇排序：做完 12 輪之後。} 前 13 個數 12, 25, \dots, 87, 89 正好是全部資料最小的 13 個，而且
已經排好、在最終位置上；後面 99, 93, 89 還沒排。 選擇排序第 $k$ 輪從第 $k$ 個位置以後選最小的換到
第 $k$ 個位置，所以做完 $k$ 輪時前 $k$ 個是最小的 $k$ 個（資結書第~14.2~節）。 照原數列實際做，
第 12 輪結束時就是這個數列（第 13 輪要放的 89 剛好已經在第 13 個位置，做完第 13 輪也不變）。
不是泡沫排序：泡沫排序是把\textbf{最大的}依序固定在\textbf{後面}，這裡後面 99, 93, 89 反而還沒排好。

\textbf{㈣ 泡沫排序：做完 3 輪之後。} 最後三個 89, 93, 99 是最大的三個（99、93 和兩個 89 之一），已經
排好在最終位置，前面 13 個還很亂。 泡沫排序每一輪把目前範圍內最大的往後推到底，3 輪後最後 3 個就定位
（資結書第~14.2~節）。 照原數列實際做 3 輪：
\begin{itemize}
\item 第 1 輪：99 推到最後。
\item 第 2 輪：93 推到倒數第二。
\item 第 3 輪：89 推到倒數第三。
\end{itemize}
結果是 32, 75, 75, 81, 89, 25, 78, 54, 75, 87, 12, 75, 28, 89, 93, 99，和數列一樣。 前面的部分也看得出
泡沫排序的痕跡：小的數（如 12、25、28）每輪只往前移一格，還留在後半段。
\end{pastsoln}
